\documentclass[10pt]{article}
\usepackage[T1]{fontenc}
\usepackage{lmodern,amsmath,amssymb,amsthm,microtype}
\usepackage[margin=0.8in]{geometry}
\usepackage[colorlinks=true,urlcolor=blue]{hyperref}
\newtheorem{theorem}{Theorem}
\newtheorem{corollary}{Corollary}
\title{Four Formulas for Even-Length Binary Palindromes Divided by Three}
\author{Qichao Wang\\{\small Hebei University of Technology}\\[0.3em]
Daoyu Dong\\{\small University of Electronic Science and Technology of China}}
\date{October 26, 2026}
\hypersetup{pdftitle={Four Formulas for OEIS A048702},pdfauthor={Qichao Wang and Daoyu Dong}}
\begin{document}
\maketitle
\begin{abstract}
OEIS A048702 lists even-length binary palindromes divided by three.
We prove both conjectured formulas and both formulas marked unproved
in its Maple section. A single carry calculation establishes a more
general identity for even-length palindromes in every integer base.
All statements hold at every index, including the changes of length.
\end{abstract}

\section{A general base identity}
Fix an integer base $b\ge2$. For $n\ge1$, write
$n=\sum_{j=0}^{L}d_jb^j$, where $d_L\ne0$ and
$L=\lfloor\log_b n\rfloor$, and put
\[
 P_b(n)=n b^{L+1}+\sum_{j=0}^{L}d_jb^{L-j},\qquad
 A_b(n)=\frac{P_b(n)}{b+1},\qquad A_b(0)=0.
\]
The integer $P_b(n)$ concatenates the base-$b$ expansion of $n$
and its reversal. Each symmetric pair of digits contributes a multiple
of $b+1$, since the two corresponding exponents have opposite parity.
Thus $A_b(n)$ is an integer. Within each fixed length the first half
orders the palindromes, and increasing its length increases the whole
number. Consequently $P_b(1),P_b(2),\ldots$ enumerate all positive
even-length base-$b$ palindromes in increasing order. In particular,
$A_2(n)$ is precisely A048702, with the initial zero at index zero.

For positive $i$, let $v_b(i)$ be the greatest $t$ with $b^t\mid i$;
the base need not be prime.
\begin{theorem}\label{thm:general}
For $n\ge1$ and $L=\lfloor\log_b n\rfloor$,
\begin{align}
 A_b(n)&=b^L\sum_{i=1}^{n}b^{-v_b(i)},\label{eq:sum}\\
 &=n b^L-(b-1)\sum_{j=1}^{L}b^{L-j}
                 \left\lfloor\frac{n}{b^j}\right\rfloor.\label{eq:floor}
\end{align}
The same sequence satisfies, with $t=v_b(n)$,
\begin{equation}\label{eq:rec}
 A_b(n)=
 \begin{cases}
 bA_b(n-1)+1,&n=b^L,\\
 A_b(n-1)+b^{L-t},&n\ne b^L.
 \end{cases}
\end{equation}
\end{theorem}
\begin{proof}
Suppose first that $n$ is not a power of $b$. The two prefixes $n-1,n$
have the same length $L+1$. In passing from $n-1$ to $n$, exactly
$t=v_b(n)$ trailing digits equal to $b-1$ become zero, and the preceding
digit increases by one. The change in the reversed half is
\[
 b^{L-t}-(b-1)\sum_{j=0}^{t-1}b^{L-j}
       =(b+1)b^{L-t}-b^{L+1}.
\]
The first half increases by $b^{L+1}$, proving the second branch of
\eqref{eq:rec}. At $n=b^L$, direct concatenation gives
\[
 P_b(n)=b^{2L+1}+1,\qquad P_b(n-1)=b^{2L}-1,
\]
where the second formula also gives zero when $L=0$. These identities
prove the first branch.

Let $F(n)$ denote the right side of \eqref{eq:sum}, and set $F(0)=0$.
Within a length block its difference is $b^{L-v_b(n)}$. At $n=b^L$
the scale factor changes from $b^{L-1}$ to $b^L$, and the newly added
summand contributes one; hence $F(n)=bF(n-1)+1$. Thus $F$ satisfies
the same recurrence and initial value as $A_b$.
Finally, for every $t\ge0$,
\[
 b^{-t}=1-(b-1)\sum_{j=1}^{t}b^{-j}.
\]
Summing this identity over $t=v_b(i)$, and counting the multiples of
$b^j$ up to $n$, proves \eqref{eq:floor}.
\end{proof}

\section{The four OEIS formulas}
The two conjectures in the formula section of A048702 are exactly
\eqref{eq:sum} and \eqref{eq:floor} at $b=2$. They therefore hold for
every $n\ge1$. The first Maple formula defines
$d=2^{\lfloor\log_2 n\rfloor-v_2(n)}$. Here $d=1$ if and only if
$n$ is a power of two: if $v_2(n)=L$, then $2^L\le n<2^{L+1}$
forces $n=2^L$. Its two branches are consequently exactly
\eqref{eq:rec}.

\begin{corollary}\label{cor:maple}
The second Maple formula also holds for every $n\ge1$:
\begin{equation}\label{eq:maple}
 A_2(n)=\frac{2}{3}(4^L-1)+1+
 \sum_{j=0}^{L}2^{L-j}
 \left\lfloor\frac{n-2^L+2^j}{2^{j+1}}\right\rfloor,
 \qquad L=\lfloor\log_2 n\rfloor.
\end{equation}
\end{corollary}
\begin{proof}
At the beginning of the block,
$A_2(2^L)=(2^{2L+1}+1)/3=2(4^L-1)/3+1$.
Write $r=n-2^L$, so $0\le r<2^L$.
For $0\le j<L$, the number of integers $i$ in
$2^L<i\le2^L+r$ with $v_2(i)=j$ is
$\lfloor(r+2^j)/2^{j+1}\rfloor$.
Each contributes $2^{L-j}$ by \eqref{eq:rec}. The $j=L$ summand
in \eqref{eq:maple} is zero, so it can be included unchanged.
\end{proof}

\section{Reproducible checks and scope}
The proof settles all four statements marked conjectured or unproved
in A048702 as retrieved on October 26, 2026. The direct palindrome
definition and its divisibility by three were already recorded there;
the new contribution is the proof of the four alternative formulas.
The general-base theorem is an additional consequence.
The accompanying \texttt{verify.py} uses exact rational arithmetic to
check the weighted sum, floor sum, carry recurrence and palindrome
definition for $1\le n\le20000$ in each base $2,\ldots,12$.
It checks \eqref{eq:maple} in base two, compares all 1001 published
b-file entries, and independently scans all integers below $2^{18}$
to verify the palindrome enumeration. These finite checks supplement
the all-index proofs.

\begin{thebibliography}{9}
\bibitem{oeis} The OEIS Foundation Inc.,
\href{https://oeis.org/A048702}{A048702} and
\href{https://oeis.org/A048701}{A048701}, accessed October 26, 2026.
\end{thebibliography}
\end{document}
